% Auto-generated by the bl-book typesetter (Bernstein Logic literate rule
% catalogue) from the mechanized rule sources, revision daaf663.
% Do not hand-edit: regenerate with `bl-book --emit-fragment` and re-copy.
% Lean statements are extracted from the sources at emission and hash-bound;
% verifier receipts were consulted at emission (their in-text badges
% record checker-layer coverage of the bound builds, not a live
% re-check by the host build).
% This file: packages + macros. Load from the host PREAMBLE, before
% hyperref. Companion body: qtsl-rule-catalogue.tex.

\usepackage{fontspec}
\usepackage{amsmath,amssymb,amsthm}
\usepackage{mathtools}
\usepackage{stmaryrd}%
\usepackage{mathpartir}
\usepackage{longtable}
\usepackage{booktabs}
\usepackage{array}
\usepackage{xcolor}
\usepackage{fancyvrb}

\directlua{luaotfload.add_fallback("bookmonofb",{
  "file:FreeMono.otf:mode=node;",
  "file:DejaVuSans.ttf:mode=node;",
  "file:FreeSerif.ttf:mode=node;",
  "file:STIX-Regular.otf:mode=node;",
  "file:Asana-Math.otf:mode=node;",
})}
\setmonofont{DejaVu Sans Mono}[Scale=MatchLowercase,RawFeature={fallback=bookmonofb}]


\definecolor{DiffRed}{RGB}{176,32,32}
\definecolor{DiffBlue}{RGB}{22,80,190}
\definecolor{MarkGreen}{RGB}{22,122,58}
\definecolor{NoteGray}{RGB}{92,92,92}
\definecolor{BadgeGreen}{RGB}{22,122,58}
\definecolor{BadgeRed}{RGB}{176,32,32}
\definecolor{BadgeAmber}{RGB}{176,108,12}
\definecolor{BadgeBlue}{RGB}{22,80,190}
\definecolor{BadgeGray}{RGB}{110,110,110}

\newcommand{\Q}{\mathbb{Q}}
\newcommand{\R}{\mathbb{R}}
\newcommand{\N}{\mathbb{N}}
\newcommand{\B}{\mathbb{B}}
\newcommand{\ReLU}{\operatorname{relu}}
\newcommand{\Cell}{\mathsf{cell}}
\newcommand{\Tensor}{\mathsf{Tensor}}
\newcommand{\sub}[2]{[#1 := #2]}
\renewcommand{\sp}{\mathsf{sp}}
\newcommand{\True}{\top}
\newcommand{\False}{\bot}
\newcommand{\rulename}[1]{\textsc{#1}}
\newcommand{\bshape}{\mathrm{bshape}}
\newcommand{\wires}{\mathrm{wires}}

\newcounter{rulesec}
\newcommand{\ruleblock}[1]{%
  \par\bigskip{\large\bfseries #1}\par\smallskip}
\newcommand{\sourcenote}[1]{{\footnotesize\textcolor{NoteGray}{#1}}\par}
\newcommand{\cataloguenote}[1]{%
  \par\smallskip{\footnotesize\textcolor{NoteGray}{\textbf{notes}\enspace #1}}\par}

\newcommand{\paradeleader}{\leaders\hbox to 0.72em{\hss.\hss}\hfill}
\newcommand{\paradeentry}[3]{%
  \par\noindent
  \sbox0{$\displaystyle #1$}%
  \ifdim\wd0>0.60\linewidth
    \usebox0\par\nopagebreak
    \hbox to \linewidth{\hfill\small\rulename{#2}\enspace\paradeleader\enspace\pageref{#3}}%
  \else
    \hbox to \linewidth{$\vcenter{\box0}$\quad{\small\rulename{#2}}\enspace\paradeleader\enspace{\small\pageref{#3}}}%
  \fi\par\medskip}
\newcommand{\paradeblock}[3]{%
  \par\noindent\begin{minipage}{\linewidth}#1\end{minipage}\par\nopagebreak
  \hbox to \linewidth{\hfill\small\rulename{#2}\enspace\paradeleader\enspace\pageref{#3}}\par\medskip}
\newcommand{\paradesep}{\par\medskip\noindent\hrule height 0.4pt\medskip}

\newcommand{\badge}[2]{%
  \fcolorbox{#1}{#1!10}{\textcolor{#1!75!black}{\sffamily\scriptsize\bfseries #2}}}

\newcommand{\declhead}[3]{%
  \par\smallskip{\small\texttt{#1}\,\texttt{.}\,\texttt{\bfseries #2}}\hfill #3\par\nopagebreak}
\newcommand{\holhead}[1]{%
  \par\smallskip{\small\texttt{\bfseries #1}}\par\nopagebreak}
\DefineVerbatimEnvironment{leansrc}{Verbatim}%
  {fontsize=\small,frame=leftline,rulecolor=\color{DiffBlue},framesep=3mm,xleftmargin=3mm}
\DefineVerbatimEnvironment{holsrc}{Verbatim}%
  {fontsize=\small,frame=leftline,rulecolor=\color{MarkGreen},framesep=3mm,xleftmargin=3mm}

\definecolor{ApiTsBlue}{RGB}{38,110,158}
\definecolor{ApiRustClay}{RGB}{164,84,36}
\newcommand{\apilang}[2]{%
  \par\medskip{\bfseries #1}\enspace{\small\texttt{#2}}\par\nopagebreak\smallskip}
\newcommand{\apiprov}[1]{{\footnotesize\textcolor{NoteGray}{#1}}\par\nopagebreak}
\DefineVerbatimEnvironment{tssrc}{Verbatim}%
  {fontsize=\footnotesize,frame=leftline,rulecolor=\color{ApiTsBlue},framesep=3mm,xleftmargin=3mm}
\DefineVerbatimEnvironment{rustsrc}{Verbatim}%
  {fontsize=\footnotesize,frame=leftline,rulecolor=\color{ApiRustClay},framesep=3mm,xleftmargin=3mm}

\newcommand{\ruleflavor}[1]{\par\smallskip #1\par}

\newenvironment{receipttable}
  {\par\scriptsize\setlength{\tabcolsep}{4pt}\begin{longtable}{@{}llll@{}}
   \textbf{receipt} & \textbf{minted} & \textbf{layers} & \textbf{status}\\
   \midrule}
  {\end{longtable}\par}
\newcommand{\receiptrow}[4]{\texttt{#1} & #2 & \texttt{#3} & #4\\}
\newcommand{\layerline}[1]{\par\smallskip{\footnotesize\textbf{gate-layer coverage of the bound build:}\par #1}\par}
\newcommand{\layeritem}[3]{{\footnotesize\quad\texttt{#1}\ \textcolor{NoteGray}{(#2)}\enspace #3}\par}

\newcommand{\absentbox}[2]{%
  \par\smallskip
  \fcolorbox{BadgeAmber}{BadgeAmber!6}{\begin{minipage}{\dimexpr\linewidth-2\fboxsep-2\fboxrule}
  {\sffamily\bfseries\textcolor{BadgeAmber!80!black}{#1}}\par\smallskip
  {\small #2}
  \end{minipage}}\par}

\usepackage{aliascnt}
\newtheorem{blemma}{Lemma}[rulesec]
\renewcommand{\theblemma}{R\arabic{rulesec}.\arabic{blemma}}
\newcommand{\blemmaautorefname}{Lemma}
\newcommand{\newaliasedbthm}[2]{%
  \newaliascnt{#1}{blemma}%
  \newtheorem{#1}[#1]{#2}%
  \aliascntresetthe{#1}%
  \expandafter\def\csname #1autorefname\endcsname{#2}%
}
\newaliasedbthm{bprop}{Proposition}
\theoremstyle{remark}
\newaliasedbthm{bremark}{Remark}
\theoremstyle{definition}
\newaliasedbthm{brecipe}{Recipe}
\theoremstyle{plain}
\newaliasedbthm{bthm}{Theorem}
\newenvironment{anchors}
  {\par\smallskip\footnotesize\begin{longtable}{@{}p{0.42\linewidth}p{0.52\linewidth}@{}}
   \textbf{proof step} & \textbf{formal anchor}\\\midrule}
  {\end{longtable}\par}

\newcommand{\dok}{\textcolor{MarkGreen}{\ensuremath{\checkmark}}}
\newcommand{\dcond}[1][hyp]{\textcolor{BadgeAmber}{\ensuremath{\checkmark}}\textsuperscript{\textcolor{BadgeAmber}{\sffamily\bfseries #1}}}
\newcommand{\dnone}{\textcolor{BadgeGray}{---}}
\newcommand{\damber}{\textcolor{BadgeAmber}{\ensuremath{\sim}}}
\newcommand{\dambercond}[1][hyp]{\textcolor{BadgeAmber}{\ensuremath{\sim}}\textsuperscript{\textcolor{BadgeAmber}{\sffamily\bfseries #1}}}
\newcommand{\dchanged}{\textcolor{BadgeAmber}{\ensuremath{\times}}}
\newcommand{\dredx}{\textcolor{BadgeRed}{\ensuremath{\times}}}
\newcommand{\dscope}{\textcolor{BadgeBlue}{\ensuremath{\varnothing}}}
\newcommand{\ddoc}{\textcolor{BadgeBlue}{\sffamily\scriptsize doc}}
\newcommand{\lcur}[1]{\textcolor{MarkGreen}{\bfseries\footnotesize #1}}
\newcommand{\lstale}[1]{\textcolor{BadgeAmber}{\footnotesize #1}}
\newcommand{\lnone}[1]{\textcolor{BadgeGray!50}{\footnotesize #1}}
\newcommand{\dashboardhead}[1]{%
  \ruleblock{Coverage dashboard}
  \sourcenote{#1}
  \sourcenote{Legend --- one row per catalogue rule (every rules-page row
  carrying a tier), ordered by tier class: core rules first, then
  admissible, then alias, then the edge-discipline glue; within a class,
  catalogue order. The small gray mark after each rule names its source
  page (BL = the Bernstein Logic pages --- the core roster and the
  composition rules, which are Bernstein Logic outright; QTSL = the QTSL
  rule roster). Lean: \dok\ statement bound and fresh; an
  \emph{amber} checkmark marks statements bound and fresh but warranting
  the rule \emph{conditionally} --- where the condition is a single named
  profile or fragment, the superscript says which (\dcond[DTMC]\ = proven
  at the discrete-kernel profile; the model-class names are the System
  chapter's), and the generic \dcond\ marks a warrant resting on a bag of
  named side conditions (the named hypotheses live on the rule's own
  catalogue page); \damber\ module
  moved since bind (\dambercond\ likewise, for a conditional-warrant row
  --- flag and label survive staleness), \dchanged\ statement changed since
  bind (stale --- the
  page badges say where), --- no Lean binding. Corr: \dok\ a
  reviewer-certified correspondence dictionary (how the Lean statements
  interpret the displayed rule) is current against the displayed cell, the
  statement texts, and its own prose; \damber\ a dictionary exists but its
  certification predates a change on some side; --- no dictionary yet. The
  correspondence claim is reviewer-grade, not machine-checked --- its
  machine-checked successor supersedes it per rule as soundness statements
  over rules-as-data land. Receipt layers
  (\texttt{b}uild/\texttt{r}eplay/\texttt{h}ardened/\texttt{d}ecorrelated):
  \lcur{green} = a current receipt covers it, \lstale{amber} = stale
  receipts only, \lnone{grey} = none. HOL4: \dok\ ported, \dscope\ out of
  scope by design, --- absent; HOL rcpt: \ddoc\ = documentary citation (no
  receipt store yet). (An Isabelle column is hidden while the Isabelle
  development is empty; it returns with its first entry.) Amber throughout
  means stale or absent; red is reserved for an outright verification
  failure. Empty cells are deliberate: the table records, rule by rule,
  what the assurance edifice does not yet cover.}
  \par\small\begin{longtable}{@{}lp{16mm}ccccc c@{}}
  \textbf{rule} & \textbf{tier} & \textbf{Lean stmt} & \textbf{corr} & \textbf{receipts} & \textbf{HOL stmt} & \textbf{HOL rcpt} & \textbf{proof}\\
  \midrule}
\newcommand{\dashboardfoot}{\end{longtable}\par}
\newcommand{\dashpage}[1]{\multicolumn{8}{@{}l@{}}{\footnotesize\textcolor{NoteGray}{page source: #1}}\\[1pt]}%
\newcommand{\dashsrc}[1]{{\scriptsize\textcolor{NoteGray}{#1}}}
\newcommand{\dashrow}[8]{#1 & \footnotesize #2 & #3 & #4 & #5 & #6 & #7 & #8\\}


\newcommand{\blcatrulebreak}{\clearpage}
\newcommand{\blcatsection}[1]{%
  \clearpage\section{#1}%
  \gdef\blcatrulebreak{\par\bigskip\gdef\blcatrulebreak{\clearpage}}}
\newcommand{\rulesection}[3]{%
  \blcatrulebreak
  \refstepcounter{rulesec}%
  \pdfbookmark[2]{Rule #1.\ #2}{blcat-rule-#1}%
  {\LARGE\bfseries Rule #1.\ \rulename{#2}}\enspace
  {\large\textcolor{NoteGray}{(#3)}}\par
  \medskip\hrule height 0.8pt\medskip}
